\documentclass[10pt,letterpaper]{article}
\usepackage{cornell-notes}

\title{Chapter 8: Network Security}
\author{}
\date{}

\setDocTitle{Chapter 8: Network Security}
\setDocSubtitle{Computer Networks}
\setDocVersion{v1.0}
\setDocStatus{Study Companion}
\setDocOwner{Computer Networks Cornell Notes}
\setDocAuthor{}
\setDocDate{\today}
\setDocSummary{Cornell-format study and review notes for Chapter 8: Network Security.}

\setCornellCollection{Computer Networks Cornell Notes}
\setCornellUnitType{Chapter}
\setCornellUnitNumber{8}
\setCornellUnitTitle{Network Security}
\setCornellCompanionLabel{Study and review companion}

\hypersetup{pdftitle={Chapter 8: Network Security},pdfsubject={Cornell notes},pdfauthor={}}

\begin{document}
\maketitle
\notesection{Chapter overview}
\begin{overviewbox}
\textbf{Learning objectives}\begin{itemize}
  \item Distinguish confidentiality, integrity, authentication, and nonrepudiation goals.
  \item Explain symmetric encryption, public-key operations, message digests, and digital signatures.
  \item Trace certificate and public-key-infrastructure trust decisions.
  \item Describe IPsec, firewalls, VPNs, wireless security, authentication protocols, and protected mail/Web channels.
  \item Recognize deprecated mechanisms and the limits of cryptography within broader systems.
\end{itemize}
\textbf{Prerequisite concepts}\quad Modular arithmetic, probability, protocol layering, IP, DNS, electronic mail, and HTTP.\par
\textbf{Key themes}\quad Threat models; secrecy; integrity; identity; freshness; key distribution; trust chains; defense in depth; privacy and policy.\par
\textbf{Named mechanisms and standards}\quad Substitution and transposition, one-time pad, DES, AES, block modes, RSA, ElGamal-like public-key methods, SHA-family digests, certificates, X.509, PKI, IPsec, firewalls, VPNs, Diffie-Hellman, Kerberos, PGP, S/MIME, SSL, and secure naming.
\end{overviewbox}

\notesection{Cornell cue-and-note record}

\begin{longtable}{@{}L{0.245\textwidth}|L{0.695\textwidth}@{}}
\rowcolor{CNavy}\color{white}\bfseries Cue / recall prompt &
\color{white}\bfseries Notes \\
\endfirsthead
\rowcolor{CNavy}\color{white}\bfseries Cue / recall prompt &
\color{white}\bfseries Notes (continued) \\
\endhead
\hline
\endfoot
\cellcolor{CCue}\textbf{8.1 Cryptography} & \begin{minipage}[t]{\linewidth}\raggedright \begin{itemize}\item \textbf{8.1.1 Introduction} A security design begins with assets, adversary capabilities, and required properties. Encryption transforms plaintext under a key; sound security depends on key secrecy and well-analyzed algorithms, not secret algorithm design.\item \textbf{8.1.2 Substitution ciphers} A substitution maps symbols to different symbols. Simple substitutions preserve language statistics and are vulnerable to frequency analysis.\item \textbf{8.1.3 Transposition ciphers} A transposition permutes symbol positions without changing the symbols. Patterns and insufficient key space permit analysis.\item \textbf{8.1.4 One-time pads} A truly random key as long as the message, used once and kept secret, provides perfect secrecy under the stated assumptions. Key generation, distribution, storage, and one-time use make it difficult to deploy.\item \textbf{8.1.5 Fundamental principles} Design the system assuming the adversary knows the algorithm, and minimize unnecessary information or access. Security should reside in keys and carefully stated assumptions.\end{itemize}\end{minipage} \\ \hline
\cellcolor{CCue}\textbf{8.2 Symmetric-key algorithms} & \begin{minipage}[t]{\linewidth}\raggedright \begin{itemize}\item \textbf{8.2.1 DES} DES is a Feistel block cipher with a key size too small for contemporary protection. Its structure is historically important, but brute-force feasibility makes single DES deprecated.\item \textbf{8.2.2 AES} AES transforms a fixed-size state through substitution, permutation, mixing, and round-key addition. Key expansion derives round keys; the design supports multiple key lengths.\item \textbf{8.2.3 Cipher modes} Modes determine how a block cipher handles multiple blocks. ECB exposes repeated patterns; chaining and feedback modes use state or counters and require correct initialization and integrity design.\item \textbf{8.2.4 Other ciphers} Other block and stream designs illustrate different performance and implementation tradeoffs; algorithm selection must follow current analysis rather than novelty.\item \textbf{8.2.5 Cryptanalysis} Attacks may exploit structure, statistics, related inputs, implementation leakage, or inadequate key size. Resistance claims are conditional on the attack model.\end{itemize}\end{minipage} \\ \hline
\cellcolor{CCue}\textbf{8.3 Public-key algorithms} & \begin{minipage}[t]{\linewidth}\raggedright Public-key cryptography separates public verification/encryption material from private signing/decryption material but is computationally heavier than symmetric encryption.\par\begin{itemize}\item \textbf{8.3.1 RSA} RSA uses a public exponent and modulus for public operations and a private exponent for private operations. Security relies on the difficulty associated with factoring the modulus; safe use requires padding and validated parameters, not raw textbook exponentiation.\item \textbf{8.3.2 Other public-key algorithms} Discrete-logarithm systems and elliptic-curve variants provide encryption, agreement, or signature capabilities with different parameter and performance profiles.\end{itemize}\end{minipage} \\ \hline
\cellcolor{CCue}\textbf{8.4 Digital signatures} & \begin{minipage}[t]{\linewidth}\raggedright \begin{itemize}\item \textbf{8.4.1 Symmetric-key signatures} A trusted intermediary or message-authentication construction can authenticate messages between parties sharing secrets, but shared knowledge complicates third-party proof.\item \textbf{8.4.2 Public-key signatures} A signer applies a private-key operation to a defined message representation; verifiers use the public key. The signature provides origin and integrity evidence only if key ownership and implementation are trustworthy.\item \textbf{8.4.3 Message digests} A cryptographic hash compresses arbitrary input to a fixed-size digest and should resist preimage, second-preimage, and collision attacks. Signatures normally cover a digest rather than an entire long message.\item \textbf{8.4.4 Birthday attack} Collision search benefits from the birthday effect: an $n$-bit digest offers roughly $2^{n/2}$ generic collision work, so collision security is about half the output length in bits.\end{itemize}\end{minipage} \\ \hline
\cellcolor{CCue}\textbf{8.5 Management of public keys} & \begin{minipage}[t]{\linewidth}\raggedright \begin{itemize}\item \textbf{8.5.1 Certificates} A certificate binds a subject name or identity to a public key and related constraints under an issuer's signature.\item \textbf{8.5.2 X.509} X.509 defines certificate fields and signed structures used in hierarchical authentication. Names, validity, algorithms, extensions, and issuer/subject relationships must all be checked.\item \textbf{8.5.3 Public-key infrastructures} A PKI manages issuance, chains, revocation information, policies, and trust anchors. Validation establishes a path to an accepted anchor; it does not prove that the endpoint itself is benign.\end{itemize}\end{minipage} \\ \hline
\cellcolor{CCue}\textbf{8.6 Communication security} & \begin{minipage}[t]{\linewidth}\raggedright \begin{itemize}\item \textbf{8.6.1 IPsec} IPsec applies security associations and packet processing at the network layer. Authentication Header protects integrity for covered fields; Encapsulating Security Payload can provide confidentiality plus integrity, in transport or tunnel mode.\item \textbf{8.6.2 Firewalls} Packet filters, stateful inspection, and application gateways restrict traffic across a boundary. A firewall enforces policy at observed paths but cannot correct vulnerable allowed services or traffic that bypasses it.\item \textbf{8.6.3 Virtual private networks} A VPN uses authenticated encryption or related protection to create a logical private path over an untrusted network, commonly through tunnels between hosts or gateways.\item \textbf{8.6.4 Wireless security} Radio is openly receivable, so link access needs authentication, key establishment, confidentiality, integrity, and replay defense. Weak protocols can fail even when encryption is present.\end{itemize}\end{minipage} \\ \hline
\cellcolor{CCue}\textbf{8.7 Authentication protocols} & \begin{minipage}[t]{\linewidth}\raggedright \begin{itemize}\item \textbf{8.7.1 Shared-secret authentication} A verifier challenges a claimant with fresh data; the claimant proves key knowledge without sending the secret. Nonces and direction-specific messages resist replay and reflection.\item \textbf{8.7.2 Diffie-Hellman} Parties exchange public values and derive a shared secret without previously sharing it. Unauthenticated exchange is vulnerable to an active intermediary, so identity binding is required.\item \textbf{8.7.3 Key distribution center} A trusted KDC shares a long-term key with each principal and distributes fresh session keys, reducing pairwise-key storage at the cost of central trust and availability.\item \textbf{8.7.4 Kerberos} Kerberos uses a KDC, tickets, authenticators, timestamps, and session keys to support single sign-on without sending passwords to each service.\item \textbf{8.7.5 Public-key authentication} Challenges, signatures, certificates, and ephemeral key agreement can authenticate peers and establish session keys, provided validation and freshness checks are correct.\end{itemize}\end{minipage} \\ \hline
\cellcolor{CCue}\textbf{8.8 Email security} & \begin{minipage}[t]{\linewidth}\raggedright \begin{itemize}\item \textbf{8.8.1 PGP} PGP combines content encryption, digital signatures, compression, textual encoding, and a user-oriented public-key trust model.\item \textbf{8.8.2 S/MIME} S/MIME protects MIME entities using certificate-based signing and encryption integrated with Internet mail formats.\end{itemize}\end{minipage} \\ \hline
\cellcolor{CCue}\textbf{8.9 Web security} & \begin{minipage}[t]{\linewidth}\raggedright \begin{itemize}\item \textbf{8.9.1 Threats} Threats include eavesdropping, tampering, impersonation, malicious input or code, tracking, and attacks on clients or servers. A protected channel does not make application logic safe.\item \textbf{8.9.2 Secure naming} A user must bind the intended service name to the correct public key. Certificates and trustworthy name resolution/validation prevent attackers from substituting endpoints.\item \textbf{8.9.3 SSL} The presented SSL/TLS design negotiates algorithms, authenticates at least the server in the common case, establishes keys, and protects records for applications such as HTTP. SSL terminology and older versions are legacy; deprecated versions must not be selected for new systems.\item \textbf{8.9.4 Mobile code security} Downloaded code is constrained through authentication, sandboxing, type safety, permissions, and runtime checks. Each mechanism limits a class of behavior but can fail through implementation flaws or excessive authority.\end{itemize}\end{minipage} \\ \hline
\cellcolor{CCue}\textbf{8.10 Social issues} & \begin{minipage}[t]{\linewidth}\raggedright \begin{itemize}\item \textbf{8.10.1 Privacy} Networks enable collection, correlation, retention, and disclosure of personal activity. Encryption can protect transit, but policy and endpoint behavior govern much of the remaining risk.\item \textbf{8.10.2 Freedom of speech} Filtering, blocking, anonymity, jurisdiction, and intermediary control create technical and social tensions; technical feasibility does not settle legitimate policy.\item \textbf{8.10.3 Copyright} Digital copying and distribution challenge enforcement and fair use. Watermarking and access controls are technical tools with limits, not complete resolutions of legal conflict.\end{itemize}\end{minipage} \\ \hline
\cellcolor{CCue}\textbf{8.11 Summary} & \begin{minipage}[t]{\linewidth}\raggedright Security composes threat models, cryptographic primitives, key management, authentication, protected channels, and operational policy. Weakness at any layer or endpoint can defeat strong algorithms.\end{minipage} \\ \hline
\end{longtable}

\begin{legacybox}Single DES, ECB for general data protection, unauthenticated Diffie-Hellman, weak wireless schemes, and obsolete SSL versions are historical or insecure mechanisms. They are included only to explain design evolution and failure modes, not as recommendations for new systems.\end{legacybox}

\notesection{Terminology and notation}

\begin{longtable}{@{}L{0.245\textwidth}|L{0.695\textwidth}@{}}
\rowcolor{CNavy}\color{white}\bfseries Cue / recall prompt &
\color{white}\bfseries Notes \\
\endfirsthead
\rowcolor{CNavy}\color{white}\bfseries Cue / recall prompt &
\color{white}\bfseries Notes (continued) \\
\endhead
\hline
\endfoot
\cellcolor{CCue}\textbf{Confidentiality} & \begin{minipage}[t]{\linewidth}\raggedright Protection against unauthorized disclosure of information.\end{minipage} \\ \hline
\cellcolor{CCue}\textbf{Integrity} & \begin{minipage}[t]{\linewidth}\raggedright Detection or prevention of unauthorized modification.\end{minipage} \\ \hline
\cellcolor{CCue}\textbf{Authentication} & \begin{minipage}[t]{\linewidth}\raggedright Evidence that an entity or message origin has the claimed identity.\end{minipage} \\ \hline
\cellcolor{CCue}\textbf{Authorization} & \begin{minipage}[t]{\linewidth}\raggedright Decision about what an authenticated principal may do.\end{minipage} \\ \hline
\cellcolor{CCue}\textbf{Nonce} & \begin{minipage}[t]{\linewidth}\raggedright Fresh value used once to prevent replay of an old protocol message.\end{minipage} \\ \hline
\cellcolor{CCue}\textbf{MAC} & \begin{minipage}[t]{\linewidth}\raggedright Message authentication code computed with a shared secret; here MAC does not mean medium access control.\end{minipage} \\ \hline
\cellcolor{CCue}\textbf{Digital signature} & \begin{minipage}[t]{\linewidth}\raggedright Private-key-generated evidence verified with a public key over a defined message representation.\end{minipage} \\ \hline
\cellcolor{CCue}\textbf{Certificate} & \begin{minipage}[t]{\linewidth}\raggedright Signed binding between a public key and named subject information or constraints.\end{minipage} \\ \hline
\cellcolor{CCue}\textbf{Trust anchor} & \begin{minipage}[t]{\linewidth}\raggedright Locally accepted root of a certificate-validation path.\end{minipage} \\ \hline
\cellcolor{CCue}\textbf{Security association} & \begin{minipage}[t]{\linewidth}\raggedright Unidirectional set of algorithms, keys, identifiers, and parameters used by IPsec.\end{minipage} \\ \hline
\end{longtable}

\notesection{Comparison matrix}

\begin{longtable}{@{}L{0.313\textwidth}L{0.313\textwidth}L{0.313\textwidth}@{}}
\toprule
\textbf{Mechanism} & \textbf{Primary property} & \textbf{Critical caveat} \\ \midrule
Symmetric encryption & Confidentiality with a shared key & Secure key distribution and integrity are separate concerns \\
Cryptographic hash & Fixed-size digest & Unkeyed hashing alone does not authenticate an origin \\
Message authentication code & Integrity and shared-key origin authentication & Either key holder can generate it \\
Digital signature & Publicly verifiable origin/integrity evidence & Depends on private-key control and public-key binding \\
Certificate & Issuer-signed key/identity binding & Validation and trust do not prove benign behavior \\
VPN & Protected logical path over an untrusted network & Endpoints and authorized application traffic remain in scope \\
\bottomrule
\end{longtable}
\notesection{Chapter synthesis}

\begin{overviewbox}\textbf{Cause-and-effect chain.} An adversary model defines required properties; primitives supply encryption, hashing, authentication, and signatures; key-management systems bind and distribute keys; protocols add freshness and identities; IP, mail, and Web mechanisms apply protection; operational controls and human policy complete the system.\par
\textbf{Common confusions.} Authentication is not authorization; encryption without integrity can permit undetected modification; a hash is not a MAC; a signature does not hide content; a certificate does not make a site trustworthy; and a firewall is not a complete application defense.\par
\textbf{Derived insight.} Most security failures are composition failures: individually sound primitives can be unsafe when keys, identities, message order, errors, or user decisions are handled incorrectly.\end{overviewbox}

\notesection{Retrieval practice}

\begin{enumerate}
  \item Define an adversary model for a message sent over an untrusted network.
  \item Why do simple substitution ciphers leak information?
  \item State the assumptions required for one-time-pad perfect secrecy.
  \item Why must single DES and ECB not be recommended for new system designs?
  \item Contrast block-cipher encryption, hashing, MACs, and digital signatures.
  \item Estimate generic collision work for a 256-bit digest using the birthday bound.
  \item Trace public-key certificate validation from an endpoint certificate to a trust anchor.
  \item Compare IPsec transport and tunnel modes.
  \item What traffic and failure classes can a firewall not address by itself?
  \item Show how an active intermediary attacks unauthenticated Diffie-Hellman.
  \item How do nonces, timestamps, and sequence values each contribute freshness?
  \item Trace Kerberos ticket use from initial authentication to service access.
  \item Compare PGP and S/MIME trust organization.
  \item Why does a protected Web channel not eliminate application-layer threats?
  \item Give one privacy issue that transit encryption cannot solve.
\end{enumerate}
\begin{CornellSummaryBox}{Rapid-review Cornell summary}Security starts with assets and adversaries. Symmetric encryption protects content efficiently; public-key mechanisms support key establishment and signatures; hashes and MACs protect representations; certificates bind keys to identities. Protocols must add freshness, validation, and correct composition. Firewalls, VPNs, IPsec, protected mail, and Web channels are layers of defense, not complete guarantees.\end{CornellSummaryBox}
\end{document}
